% =====================================================================
%  TP4, Phase III, second half: results and discussion.
%  Author: FOMETHE SOBMBANANG MAXIMILIEN
%  Fragment, included by phase3_sections.tex. See that file for the
%  packages the host document must load.
% =====================================================================

\section{Results and discussion}
\label{sec:results}

\subsection{The seven baselines}
\label{sec:results-baselines}

\Cref{tab:baselines} reports the retained candidate of each family, and
\cref{fig:pr,fig:roc} show the corresponding curves on the test split.
Histogram gradient boosting leads on validation average precision with
\num{0.2719}, ahead of the RBF kernel machine at \num{0.2417} and the linear
support vector machine at \num{0.2158}. The $k$-nearest neighbours rule with
$K=11$ follows closely at \num{0.2128}, then the depth~5 tree at \num{0.2058},
Gaussian naive Bayes at \num{0.1845} and the gradient ascent logistic model at
\num{0.1833}. The first three were therefore carried into Task~3.2.

\begin{table*}[t]
  \centering
  \caption{The seven baselines on the primary seed. Each row is the candidate of
  its family with the highest validation average precision; the test columns are
  computed once, after that choice. The score type is \emph{prob.} when the
  ranking metrics come from \texttt{predict\_proba} and \emph{margin} when they
  come from \texttt{decision\_function}. Accuracy, precision, recall and F1 are
  taken at each family's default cut-off, which \cref{sec:results-operating}
  shows to be the wrong reading of this task.}
  \label{tab:baselines}
  \footnotesize\setlength{\tabcolsep}{4.5pt}
  \begin{tabular}{llcccccccc}
\toprule
\multirow{2}{*}{Model} & \multirow{2}{*}{Configuration} & \multirow{2}{*}{Score type} & Val. & \multicolumn{6}{c}{Test split} \\
\cmidrule(l){5-10}
 & & & AP & Acc. & Prec. & Rec. & F1 & AUC & AP \\
\midrule
Histogram gradient boosting & leaves 31 & prob. & \textbf{0.2719} & 0.8738 & 0.3879 & 0.0490 & 0.0870 & 0.7447 & 0.2753 \\
Kernel SVM (RBF) & $C = 1$ & margin & 0.2417 & 0.8772 & 0.0000 & 0.0000 & 0.0000 & 0.7175 & \textbf{0.2763} \\
Linear SVM (soft margin) & $C = 0.1$ & margin & 0.2158 & 0.8725 & 0.3174 & 0.0338 & 0.0610 & 0.6911 & 0.2281 \\
K-nearest neighbours & $K = 11$ & prob. & 0.2128 & 0.8740 & 0.3237 & 0.0243 & 0.0453 & 0.6732 & 0.2041 \\
Decision tree (CART) & depth 5 & prob. & 0.2058 & 0.8732 & 0.2443 & 0.0156 & 0.0294 & 0.6695 & 0.2075 \\
Gaussian Naive Bayes & $\varepsilon = 10^{-9}$ & prob. & 0.1845 & 0.6117 & 0.1923 & 0.6761 & 0.2995 & 0.6626 & 0.1898 \\
Logistic regression (SGA) & $\lambda = 10^{-5}$ & prob. & 0.1833 & 0.8522 & 0.2498 & 0.1017 & 0.1445 & 0.6606 & 0.2024 \\
\bottomrule
\end{tabular}

\end{table*}

The informative comparison is not the raw value but the lift over the no-skill
level, because average precision inherits the prevalence of the split it is
measured on, and the two held out splits do not have the same prevalence
(\num{0.1152} on validation, \num{0.1228} on test). On test, boosting reaches
\num{2.24} times the no-skill level and the RBF machine \num{2.25} times, while
the gradient ascent logistic model reaches \num{1.65}. Two conclusions follow.
The models do extract a real signal, between \num{1.6} and \num{2.3} times what
chance gives on this prevalence. And that signal is modest: a precision near
\num{0.26} at a recall near \num{0.5} (\cref{tab:operating}) means three false
alarms for every correct call, which is not a usable screening operating point.

\begin{figure}[t]
  \centering
  \includegraphics[width=\columnwidth]{fig1_pr_baselines.pdf}
  \caption{Precision-recall curves on the test split. The dotted line is the
  no-skill level, equal to the \num{0.1228} prevalence. The vertical axis is
  truncated at \num{0.70}: below a recall of about \num{0.02} the curves are
  decided by a handful of images and swing between \num{0} and \num{1}.}
  \label{fig:pr}
\end{figure}

\begin{figure}[t]
  \centering
  \includegraphics[width=\columnwidth]{fig2_roc_baselines.pdf}
  \caption{ROC curves for the same seven models and the same score vectors as
  \cref{fig:pr}. The two views do not agree on the ordering of the top two
  models, and the ROC view makes all seven look closer to one another than the
  precision-recall view does.}
  \label{fig:roc}
\end{figure}

Three observations deserve to be separated, because each is a different way the
imbalance distorts a familiar metric.

\paragraph{Accuracy ranks a constant classifier first.}
Six of the seven models sit between \num{0.8522} and \num{0.8772} test accuracy,
inside the band around the \num{0.8772} that an always-negative answer scores.
The highest accuracy of the table, \num{0.8772}, belongs to the RBF kernel
machine, and it is exactly that trivial value because at margin~\num{0} the
model makes \emph{no} positive call at all on the \num{22433} test images:
\num{0} true positives and \num{0} false positives. The model with the best test
average precision of the table is therefore also the one whose hard predictions
carry no information. Accuracy appears in \cref{tab:baselines} for completeness
and is used nowhere else in this work.

\paragraph{F1 at the default cut-off inverts the ranking.}
Gaussian naive Bayes has the second worst average precision (\num{0.1898} on
test) and the best default F1 (\num{0.2995}), while boosting has the best
average precision and an F1 of \num{0.0870}. The cause is visible in the number
of positive calls: naive Bayes issues \num{9681} of them on \num{22433} images,
against \num{348} for boosting and \num{0} for the RBF machine. Multiplying
\num{784} near-duplicate likelihood terms over strongly correlated pixels drives
the posterior to the ends of $[0,1]$, which leaves the default \num{0.5} cut-off
in a region where almost every image is called positive; the F1 maximising
threshold selected on validation for that model is \num{1.0000}, the boundary of
the interval, which is the signature of a saturated posterior. Naive Bayes is
not better at detecting effusion, it is differently miscalibrated. Any
conclusion drawn from F1 at the default threshold on this task would be an
artefact of calibration.

\paragraph{ROC and precision-recall disagree where it matters.}
Boosting dominates on ROC AUC (\num{0.7447} against \num{0.7175} for the RBF
machine) but the two are level on test average precision (\num{0.2753} against
\num{0.2763}). The reversal is the arithmetic of \cref{sec:protocol-metrics}:
the \num{19679} negative test images in the denominator of the false positive
rate absorb precisely the errors that the precision-recall view charges to the
model. \Cref{fig:roc} also compresses the seven models into a narrow band, which
would suggest that the choice of family hardly matters, whereas \cref{fig:pr}
separates boosting and the kernel machine from the rest over most of the recall
range. This is the concrete form, on this dataset, of the argument of Davis and
Goadrich~\cite{davis2006relationship} and of Saito and
Rehmsmeier~\cite{saito2015precision}.

One cost belongs next to the scores. The RBF machine needs \SI{146.7}{\second}
to score the \num{22433} test images against \SI{2.1}{\second} for boosting, a
factor of \num{69}, because its decision function is a sum over several thousand
support vectors in \num{784} dimensions. At equal average precision, boosting is
the better engineering choice.

\subsection{Where the decision boundary sits}
\label{sec:results-operating}

\Cref{tab:operating} reads the same seven models at a second operating point,
the threshold that maximises F1 on validation, transferred unchanged to test.
Every model improves and the degenerate cases disappear: boosting goes from
\num{0.0870} to \num{0.3639}, the RBF machine from \num{0.0000} to \num{0.3411},
the linear machine from \num{0.0610} to \num{0.3070}. Naive Bayes, first at the
default cut-off with \num{0.2995}, moves only to \num{0.3099} and falls back to
the middle of the table. The ordering by F1 at the selected threshold then
broadly follows the ordering by average precision, which is the expected
behaviour: once the threshold stops being an uncontrolled variable, the
thresholded metric measures the same underlying ranking quality.

\begin{table*}[t]
  \centering
  \caption{Test precision, recall and F1 at the default cut-off of each family
  (\num{0.5} on a posterior, \num{0} on a margin) and at the cut-off that
  maximises validation F1. The test split is never used to choose the threshold.
  Every model gains, and the two models that made almost no positive call at the
  default cut-off gain the most.}
  \label{tab:operating}
  \small
  \begin{tabular}{lcccccc}
\toprule
\multirow{2}{*}{Model} & \multicolumn{3}{c}{Default cut-off} & \multicolumn{3}{c}{Cut-off chosen on validation} \\
\cmidrule(lr){2-4}\cmidrule(l){5-7}
 & Prec. & Rec. & F1 & Prec. & Rec. & F1 \\
\midrule
Histogram gradient boosting & 0.388 & 0.049 & 0.087 & 0.264 & 0.587 & 0.364 \\
Kernel SVM (RBF) & 0.000 & 0.000 & 0.000 & 0.264 & 0.483 & 0.341 \\
Linear SVM (soft margin) & 0.317 & 0.034 & 0.061 & 0.235 & 0.442 & 0.307 \\
K-nearest neighbours & 0.324 & 0.024 & 0.045 & 0.213 & 0.570 & 0.310 \\
Decision tree (CART) & 0.244 & 0.016 & 0.029 & 0.231 & 0.490 & 0.314 \\
Gaussian Naive Bayes & 0.192 & 0.676 & 0.299 & 0.207 & 0.614 & 0.310 \\
Logistic regression (SGA) & 0.250 & 0.102 & 0.145 & 0.205 & 0.484 & 0.288 \\
\bottomrule
\end{tabular}

\end{table*}

\begin{figure}[t]
  \centering
  \includegraphics[width=\columnwidth]{fig5_threshold_sensitivity.pdf}
  \caption{F1 of the boosted model against its decision threshold, on validation
  and on test. The default \num{0.5} cut-off lies far to the right of the
  maximum, which on this prevalence sits near \num{0.07}. The two curves peak at
  nearly the same place, so the threshold transfers across splits.}
  \label{fig:threshold}
\end{figure}

\Cref{fig:threshold} shows why. For the boosted model the F1 maximum lies at a
posterior of \num{0.0691}, not at \num{0.5}, and the validation and test curves
peak at nearly the same place, so the choice transfers. The practical reading is
that a thresholded metric on an imbalanced task reports the sum of two things,
the quality of the ranking and the accident of where the library put the
cut-off, and that the second term dominates here. This is why average precision
drives every selection decision in this work, and why \cref{tab:robustness}
reports F1 at the validation selected point rather than at the default one.

\subsection{Does reducing the dimension help?}
\label{sec:results-representation}

The two families whose behaviour depends most directly on the ambient dimension
were rerun on a \num{50} component principal projection fitted on the training
subsample, retaining \num{95.3} percent of the variance
(\cref{tab:representation}). Neither improves. The $K=5$ neighbours rule falls
from \num{0.1822} to \num{0.1776} validation average precision and the RBF
machine from \num{0.2417} to \num{0.2344}. The expectation behind the ablation,
that Euclidean distances and an isotropic kernel would behave better once
\num{784} correlated pixel coordinates are compressed, is not met. The
discriminative signal lies partly in directions that principal component
analysis discards for carrying little variance, which is plausible on this
modality: variance in a chest radiograph is dominated by patient size, framing
and exposure, none of which is the costophrenic angle. The projection does buy
computation, since the RBF machine then scores the test split in
\SI{21.9}{\second} instead of \SI{146.7}{\second}, a factor of \num{6.7}, for
\num{0.008} of average precision. That trade is worth taking in deployment and
not worth taking when the point is to compare model families, so the main tables
keep the \num{784} dimensional features.

\begin{table}[t]
  \centering
  \caption{Representation check on the primary seed. The principal basis is
  fitted on the training subsample only and retains \num{95.3} percent of the
  variance.}
  \label{tab:representation}
  \footnotesize\setlength{\tabcolsep}{4pt}
  \begin{tabular}{lcccc}
\toprule
\multirow{2}{*}{Model} & \multicolumn{2}{c}{Pixels (784)} & \multicolumn{2}{c}{PCA (50)} \\
\cmidrule(lr){2-3}\cmidrule(l){4-5}
 & AUC & AP & AUC & AP \\
\midrule
KNN, $K = 5$ & 0.6336 & 0.1819 & 0.6317 & 0.1810 \\
Kernel SVM, $C = 1$ & 0.7175 & 0.2763 & 0.7086 & 0.2683 \\
\bottomrule
\end{tabular}

\end{table}
